% ==============================================================================
\chapter{Introduzione e importanza della deconvoluzione}
% ==============================================================================

\section{Introduzione al problema inverso}

Così come molti parametri non possono essere direttamente misurati \textit{in vivo}, può accadere lo stesso per molti segnali di interesse alla comprensione del sistema biologico.

\begin{itemize}
    \item \textbf{Esempio:} Il tasso di secrezione di una ghiandola deve essere inferito a partire dalla concentrazione dell'ormone nel plasma (è possibile misurare solo l'effetto e non la causa).
    \item \textbf{PROBLEMA INVERSO:} se il segnale di input del sistema non è noto, il problema inverso è un problema di stima dell'ingresso.
\end{itemize}

\noindent\textbf{Schema concettuale del sistema fisiologico (Deconvoluzione per sistemi LTI):}
\begin{center}
\begin{tabular}{ccc}
    \textbf{Ingresso non noto} & \textbf{Sistema Fisiologico} & \textbf{Uscita misurata} \\
    \textit{Production? Secretion?} & \textit{Physiological System} & \textit{Plasma Concentration} \\
    \textit{Rate of Appearance} & & \\
    $u(t)$ & $\longrightarrow \quad g(t-\tau) \quad \longrightarrow$ & $c(t)$ \\
    \multicolumn{3}{c}{\textbf{DECONVOLUZIONE (sistemi LTI)}}
\end{tabular}
\end{center}

\subsection{Importanza della deconvoluzione }

Si potrebbe argomentare che la misura di concentrazione di una sostanza nel sangue fornisce sufficienti informazioni per comprendere il sistema e quindi non è necessario risolvere in $u(t)$ l'integrale perché $c(t)$ è sufficiente:
\[
    c(t) = \int_{-\infty}^{t} g(t-\tau) u(\tau) \, d\tau
\]
Tuttavia, le concentrazioni nel plasma non riflettono soltanto il processo di secrezione/produzione a cui siamo interessati, ma risentono anche della dinamica delle sostanze (distribuzione e metabolismo nell'organismo).

\begin{center}
    \textbf{DECONVOLUZIONE: rimuovere la distorsione introdotta dalla dinamica della sostanza}
\end{center}

\noindent\textbf{Confronto tra cinetica rapida e cinetica lenta ($g(t) = e^{-\alpha t}$) }

\noindent Nell'analisi comparativa tra un sistema rapido (\textit{'Fast System'}, $\alpha = 12$) e un sistema lento (\textit{'Slow System'}, $\alpha = 2$) sottoposti al medesimo ingresso impulsivo bifasico, si osservano quattro livelli di segnale nel dominio del tempo ($t \in [0, 1]$):
\begin{figure}[h]
    \centering
    \includegraphics[width=1\textwidth]{Immagini/confrontocinetica.png}
    \label{fig:7-1}
\caption{Influence of substance kinetics on concentration time series. Left and right panels describe a 'fast' and a 'slow' system, respectively. From top to bottom panels: impulse response, input, continuous output (continuous line) and noisy samples (open circles) and measurement time series (open circles, interpolation also shown) are shown, respectively (adapted from Sparacino et al., 2004).}
\end{figure}
\iffalse
\begin{center}
\begin{tabular}{p{4.2cm}p{5.0cm}p{5.0cm}}
    \toprule
    \textbf{Grandezza} & \textbf{'Fast System' ($\alpha = 12$)} & \textbf{'Slow System' ($\alpha = 2$)} \\
    \midrule
    \textit{Impulse Response} $g(t)$ & Decadimento rapido da $1$ a $0$ ($\text{l}^{-1}$) & Decadimento lento da $1$ a $\approx 0.14$ ($\text{l}^{-1}$) \\
    \textit{Input} $u(t)$ & Due picchi a $t = 0.2$ e $t = 0.4$ ($\text{pmol/min}$) & Due picchi a $t = 0.2$ e $t = 0.4$ ($\text{pmol/min}$) \\
    \textit{Output and Noisy Samples} & Uscita continua (range $0\text{--}0.1$) e campioni rumorosi (cerchi aperti) & Uscita continua fortemente smussata/ritardata (range $0\text{--}0.04$) e campioni rumorosi \\
    \textit{Measured Time Series} & Serie temporale misurata e interpolazione ($\text{pmol/l}$, asse $x$: \textit{Time}) & Serie temporale misurata e interpolazione ($\text{pmol/l}$, asse $x$: \textit{Time}) \\
    \bottomrule
\end{tabular}
\end{center}
\fi
% ==============================================================================
\section{Formulazione matematica e criticità del problema}
% ==============================================================================

\subsection{Il problema e le assunzioni }

Risolviamo l'integrale
\[
    c(t) = \int_{-\infty}^{t} g(t-\tau) u(\tau) \, d\tau
\]
in $u(t)$, essendo noti l'uscita del sistema $c(t)$ e la sua risposta impulsiva $g(t)$.

\noindent\textbf{Assunzioni:}
\begin{enumerate}
    \item $u(t)$ è \textbf{causale}, ovvero $u(t) = 0$ per $t < 0$ (non sempre vero $\rightarrow$ secrezione ormonale spontanea). Se mi rifaccio alla fisiologia questa assunzione non ha senso, supponiamo semplicemente che la misura effettuato in zero sia il minimo delle concentrazioni che misureremo.
    \item Sistema lineare e tempo invariante, altrimenti l'integrale diventa:
    \[
        c(t) = \int_{-\infty}^{t} g(t,\tau) u(\tau) \, d\tau \qquad \text{(integrale di Fredholm)}
    \]
\end{enumerate}

\subsection{Natura mal-posta e mal-condizionata della deconvoluzione}

La deconvoluzione è un problema:
\begin{enumerate}
    \item \textbf{mal-posto (sotto-determinato)} $\rightarrow$ il problema di ricostruire l'input $u(t)$ da un numero finito di campioni dell'uscita $c(t)$ non ammette un'unica soluzione; questo peroblema può essere risolto 
    \item \textbf{mal-condizionato} $\rightarrow$ una piccola percentuale di errore nella misura produce un elevato errore nella stima dell'ingresso. (Può essere mitigato ma non può essere cancellato)
\end{enumerate}

\noindent\textbf{Esempio di deconvoluzione}
Consideriamo l'ingresso dato da (unità arbitrarie, \textit{Arbitrary Units}):
\[
    u(t) = e^{-[(t-400)/75]^2} + e^{-[(t-600)/75]^2}, \qquad 0 \le t \le 1025
\]
\begin{figure}[h]
    \centering
    \includegraphics[width=0.5\textwidth]{Immagini/esempiodiconvoluzione.png}
    \label{fig:7-2}
\end{figure}

\noindent E la risposta impulsiva del sistema (\textit{Impulse Response}, \textit{Arbitrary Units}):
\[
    g(t) = \begin{cases} 1, & t \le 250 \\ 0, & t > 250 \end{cases}
\]
Noti $u(t)$ e $g(t)$ è possibile calcolare $c(t)$.

Assumiamo invece di conoscere $N$ campioni di $c(t)$: $\{c_k\} = c(t_k)$ misurati senza errore su una griglia uniforme $\Omega_s = \{kT\}$, $k = 1, 2, \dots, N$ con $T = 25$ e $N = 41$. (Abbiamo una griglia di 1025 punti)
\begin{itemize}
    \item Il problema di ricostruire una funzione tempo-continua $u(t)$ dalla serie temporale $\{c_k\}$ può \textbf{NON} risultare in un'unica soluzione.
    \item Oltre alla soluzione vera, esisteranno altre soluzioni in grado di fornire perfettamente $\{c_k\}$ in uscita (problema mal posto).
    \item Per trattare la mal posizione del problema, gli approcci di deconvoluzione devono restringere il campo di applicazione delle funzioni in una zona in cui la soluzione è unica (\textbf{DECONVOLUZIONE DISCRETA}).
\end{itemize}

\begin{figure}[h]
    \centering
    \includegraphics[width=1\textwidth]{Immagini/esempiodiconvoluzione2.png}
    \label{fig:7-3}
\end{figure}
Da 41 valori vorrei poter stimare solo 41 punti invece che infiniti, in questo modo il problema diventerebbe quadro e quindi con un'unica soluzione.
% ==============================================================================
\section{Deconvoluzione discreta}
% ==============================================================================

\subsection{Discretizzazione dell'ingresso e dell'integrale di convoluzione}

Il segnale $u(t)$ si assume essere costante all'interno di ogni intervallo della griglia di campionamento $\Omega_s = \{t_1, t_2, \dots, t_N\}$, (costante a tratti), ovvero:

\begin{figure}[h]
    \centering
    \includegraphics[width=1\textwidth]{Immagini/deconvoluzionediscreta.png}
    \label{fig:7-4}
\end{figure}
\[
    u(t) = u_i \quad \text{per } t_{i-1} \le t \le t_i, \qquad i = 1, 2, \dots, N
\]
Sostituendo le espressioni discrete l'integrale di convoluzione diventa:
\[
    c(t_k) = \int_{0}^{t_k} g(t_k - \tau) \cdot u(\tau) \, d\tau = \sum_{i=1}^{k} u_i \cdot \int_{t_{i-1}}^{t_i} g(t_k - \tau) \, d\tau
\]
Considerando:
\[
    g_{k-i} = \int_{t_{i-1}}^{t_i} g(t_k - \tau) \, d\tau
\]
Si ottiene:
\[
    c_k = c(t_k) = \sum_{i=1}^{k} u_i \cdot g_{k-i}
\]
I cui primi termini sono:
\begin{align*}
    k = 1 &\Rightarrow c_1 = g_0 \cdot u_1 \\
    k = 2 &\Rightarrow c_2 = g_0 \cdot u_2 + g_1 \cdot u_1 \\
    k = 3 &\Rightarrow c_3 = g_0 \cdot u_3 + g_1 \cdot u_2 + g_2 \cdot u_1
\end{align*}
Questa deconvoluzione viene quindi risolta in modo discreto.\\

\noindent In notazione matriciale, ovvero, esplicitando vettori e matrice (con $g_N$ indicato nella slide per l'elemento in basso a sinistra, corrispondente a $g_{N-1}$):
\[
    c = G u \qquad \begin{bmatrix}
        c_1 \\
        c_2 \\
        c_3 \\
        \vdots \\
        c_N
    \end{bmatrix}
    =
    \begin{bmatrix}
        g_0 & 0 & 0 & \dots & 0 \\
        g_1 & g_0 & 0 & \dots & 0 \\
        g_2 & g_1 & g_0 & \dots & 0 \\
        \vdots & \vdots & \vdots & \ddots & \vdots \\
        g_{N-1} & \dots & g_2 & g_1 & g_0
    \end{bmatrix}
    \begin{bmatrix}
        u_1 \\
        u_2 \\
        \vdots \\
        u_N
    \end{bmatrix}
\]


Dove:
\begin{itemize}
    \item $c = [c_1, c_2, \dots, c_N]^T$ è il vettore $N$-dimensionale dell'uscita campionata;
    \item $u = [u_1, u_2, \dots, u_N]^T$ è l'ingresso;
    \item $G$ è una matrice triangolare inferiore $N \times N$.
\end{itemize}

\noindent Poiché $G$ è invertibile, si ottiene un'\textbf{unica soluzione}:
\[
    \hat{u} = G^{-1} c
\]
Una volta convoluto $\hat{u}$ con $g(t)$, l'uscita è perfettamente descritta.
\medskip
\noindent\textbf{Esempio deconvoluzione discreta}

\begin{figure}
    \centering
    \includegraphics[width=1\textwidth]{Immagini/esempiodeconvoluzionediscreta.png}
    \label{fig:7-5}
\end{figure}

% ==============================================================================
\section{Deconvoluzione in presenza di rumore e mal-condizionamento}
% ==============================================================================

\subsection{Modello di misura con rumore e stima ai minimi quadrati (LS)}

\noindent\textbf{Schema del sistema rumoroso }
\begin{figure}[h]
    \centering
    \includegraphics[width=1\textwidth]{Immagini/deconvoluzioneinpresenzadirumore.png}
    \label{fig:7-6}
\end{figure}

I campioni misurati in uscita sono affetti da errore di misura e questo complica il problema di deconvoluzione. Consideriamo la $k$-esima misura $z_k$:
\[
    z_k = c_k + v_k \qquad k = 1, 2, \dots, N
\]
dove $v_k$ è l'errore.

In notazione matriciale:
\[
    z = G \cdot u + v
\]
dove:
\begin{itemize}
    \item $z = [z_1, z_2, \dots, z_N]^T$ è il vettore delle misure;
    \item $v = [v_1, v_2, \dots, v_N]^T$ è il vettore degli errori di misura;
    \item $v$ è un vettore a media nulla, campioni indipendenti e matrice di covarianza nota pari a $\Sigma_v$.
\end{itemize}

La più semplice stima di $u$ da $z$ è data dalla soluzione del problema ai minimi quadrati (\textit{Least Squares}), totalmente equivalente alla $w_{rss}$:
\[
    \min_{u} \underbrace{(z - G \cdot u)^T}_{\text{residuo}} \underbrace{\Sigma_v^{-1}}_{\substack{\text{Peso}}} (z - G \cdot u)
\]

\noindent La cui soluzione è:
\[
    \hat{u}_{LS} = (G^T \Sigma_v^{-1} G)^{-1} G^T \Sigma_v^{-1} z = G^{-1} (G^T \Sigma_v^{-1})^{-1} G^T \Sigma_v^{-1} z
\]
poiché $G$ è invertibile.

In questo caso la riconvoluzione, ovvero la predizione dell'uscita $z(t)$, è:
\[
    z(t) = \int_{0}^{t} g(t-\tau) \cdot \hat{u}_{LS}(\tau) \, d\tau
\]
che descrive perfettamente i dati campionati, ovvero il vettore dei residui $r = z - G\hat{u}_{LS}$ è identicamente nullo.

\subsection{Effetti del rumore, della frequenza di campionamento e della dinamica }

\noindent\textbf{Effetti del rumore sulla stima}

La presenza del rumore nel vettore delle misure ha un effetto drammatico sulla qualità della stima. Aggiunta di rumore gaussiano ($\text{SD} = 3$) alle misure del problema precedente:
\begin{figure}[h]
    \centering
    \includegraphics[width=0.95\textwidth]{Immagini/effettorumorestima.png}
    \label{fig:7-7}
\end{figure}
\begin{itemize}
    \item Nei pannelli \textit{Input}, \textit{Impulse Response}, \textit{Ill-Posedness: Two Possible Solution} e \textit{Output with Noise-Free Samples} si ripropone il caso ideale.
    \item Nel pannello \textit{Output Noisy Samples \& Reconvolution} (asse $x$: \textit{Time}, $0\text{--}1000$; asse $y$: \textit{Arbitrary Units}, $0\text{--}200$), la riconvoluzione passa esattamente per i punti rumorosi.
    \item Nel pannello \textit{Ill-conditioning: LS Deconvolution} (asse $x$: \textit{Time}, $0\text{--}1000$; asse $y$: \textit{Arbitrary Units}, $-0.5\text{--}2$), la stima $\hat{u}_{LS}$ presenta ampie oscillazioni spurie che mascherano completamente il profilo dei due picchi originali.
\end{itemize}

\noindent\textbf{Incremento del numero di campioni}

Confronto tra passo di campionamento $T = 25$ (\textit{Ill-conditioning: LS Deconvolution} e \textit{Output Noisy Samples \& Reconvolution}) e passo di campionamento più fitto $T = 10$ (\textit{Ill-conditioning: Effect of Sampling} e \textit{Output Noisy Samples \& Reconvolution}, ordinate tra $-1$ e $2$):
\begin{figure}[h]
    \centering
    \includegraphics[width=0.8\textwidth]{Immagini/incrementocampioni.png}
    \label{fig:7-8}
\end{figure}
\begin{center}
    \textbf{Maggiore è la frequenza di campionamento più sarà difficile ricostruire il tasso di secrezione dell'ormone mediante deconvoluzione}
\end{center}

\noindent\textbf{Risposta impulsiva lenta}

Confronto a parità di passo di campionamento ($T = 25$) tra un sistema in cui $g(t)$ dura $250$ campioni (\textit{Ill-conditioning: LS Deconvolution}) e un sistema più lento in cui $g(t)$ dura $400$ campioni (\textit{Ill-conditioning: Effect of System Dynamics}):
\begin{figure}[h!]
    \centering
    \includegraphics[width=0.8\textwidth]{Immagini/rispostalenta.png}
    \label{fig:7-9}
\end{figure}
\begin{center}
    \textbf{Maggiore è il tempo di secrezione dell'ormone più sarà difficile ricostruire il tasso di secrezione dell'ormone mediante deconvoluzione}
\end{center}

% ==============================================================================
\section{Approccio non parametrico: Metodo di regolarizzazione}
% ==============================================================================

\subsection{Strategie contro il mal-condizionamento}

La deconvoluzione LS è semplice ma molto debole perché troppo sensibile al mal-condizionamento del problema. Esistono diversi approcci per limitare il mal-condizionamento:
\begin{enumerate}
    \item \textbf{Deconvoluzione parametrica} $\rightarrow$ si assume un'espressione analitica per l'input e si stimano i parametri che la descrivono (la deconvoluzione diventa un problema di stima dei parametri).
    \item \textbf{Deconvoluzione non parametrica} $\rightarrow$ non richiede la formulazione analitica a priori dell'input (approccio di regolarizzazione).
\end{enumerate}

\subsection{Formulazione del metodo di regolarizzazione}
Il metodo di regolarizzazione è un approccio non parametrico. L'idea alla base del metodo è di identificare una soluzione che non massimizzi soltanto la perfezione del fitting dei dati (come nella deconvoluzione LS), ma una soluzione che fornisca un buon fitting dei dati e simultaneamente garantisca un certo livello di \textit{«smoothness»}.
\[
\min_{\mu}{w_{rss}} \longrightarrow \min_{\mu}{(w_{rss} + \text{smooth})}
\]
Il problema di ottimizzazione da risolvere diventa:
\[
    \min_{u} (z - G \cdot u)^T \Sigma_v^{-1} (z - G \cdot u) + \gamma u^T F^T F u
\]
dove:
\begin{itemize}
    \item $\Sigma_v$ è la matrice di covarianza dell'errore di misura;
    \item $F$ è una matrice quadrata di penalizzazione ($N \times N$);
    \item $\gamma$ è un parametro reale non negativo.
\end{itemize}
Ovviamente c'è un trade-off tra fitting dei dati e smoothness della soluzione. La scelta di $\gamma$ è fondamentale per ottenere una stima soddisfacente, la smoothness è proporzionale all'errore che compieremo.

A fronte della soluzione ai minimi quadrati $\hat{u}_{LS} = (G^T \Sigma_v^{-1} G)^{-1} G^T \Sigma_v^{-1} z$, il problema regolarizzato posto è quadratico e la sua soluzione è:
\[
    \hat{u} = (G^T \Sigma_v^{-1} G + \gamma F^T F)^{-1} G^T \Sigma_v^{-1} z
\]
che dipende linearmente dal vettore dei dati $z$.
\begin{itemize}
    \item Se $\gamma = 0$, la funzione di costo e la sua soluzione coincidono con la deconvoluzione LS.
    \item Se $\gamma > 0$, la funzione di costo è composta da due termini:
    \begin{itemize}
        \item Il primo penalizza la distanza, pesata dall'inverso di $\Sigma_v$, tra la predizione del modello $Gu$ (vettore di riconvoluzione) e i dati.
        \item Il secondo contributo $u^T F^T F u$ penalizza l'irregolarità della soluzione.
    \end{itemize}
\end{itemize}

\subsection{Costruzione della matrice di penalizzazione $F$}

\noindent\textbf{Penalizzazione (1/2):} La scelta standard è di penalizzare l'energia della derivata temporale $m$-esima, con $m$ parametro intero. In genere si utilizza l'energia della derivata prima o seconda. Come si costruisce $F$?

Consideriamo il caso di campionamento uniforme. $F$ viene costruita in modo tale che $Fu$ fornisce l'$m$-esima differenza di $u$ ($m = 1, 2, \dots$):
\[
    F = D^m
\]
\[
    D = \begin{bmatrix}
        1 & 0 & 0 & \dots & 0 \\
        -1 & 1 & 0 & \dots & 0 \\
        0 & -1 & 1 & \dots & 0 \\
        \vdots & \vdots & \vdots & \ddots & \vdots \\
        0 & \dots & \dots & -1 & 1
    \end{bmatrix}
\]

\noindent\textbf{Penalizzazione (2/2):}
\begin{itemize}
    \item \textbf{Se $m = 1$:}
    \[
        F = \begin{bmatrix}
            1 & 0 & 0 & \dots & 0 \\
            -1 & 1 & 0 & \dots & 0 \\
            0 & -1 & 1 & \dots & 0 \\
            \vdots & \vdots & \vdots & \ddots & \vdots \\
            0 & \dots & \dots & -1 & 1
        \end{bmatrix}
        Du = \begin{bmatrix}
            1 & 0 & 0 & \dots & 0 \\
            -1 & 1 & 0 & \dots & 0 \\
            0 & -1 & 1 & \dots & 0 \\
            \vdots & \vdots & \vdots & \ddots & \vdots \\
            0 & \dots & \dots & -1 & 1
        \end{bmatrix}
        \begin{bmatrix}
            u_1 \\
            u_2 \\
            u_3 \\
            \vdots \\
            u_N
        \end{bmatrix}
        =
        \begin{bmatrix}
            u_1 \\
            u_2 - u_1 \\
            u_3 - u_2 \\
            \vdots \\
            u_N - u_{N-1}
        \end{bmatrix}
    \]
    \item \textbf{Se $m = 2$:}
    \[
        F = \begin{bmatrix}
            1 & 0 & 0 & \dots & 0 \\
            -2 & 1 & 0 & \dots & 0 \\
            1 & -2 & 1 & \dots & 0 \\
            \dots & 1 & -2 & \dots & 0 \\
            \vdots & \vdots & \vdots & \ddots & \vdots \\
            0 & \dots & 1 & -2 & 1
        \end{bmatrix}
    \]
    \[
        D^2 u = D(Du) = \begin{bmatrix}
            1 & 0 & 0 & \dots & 0 \\
            -1 & 1 & 0 & \dots & 0 \\
            0 & -1 & 1 & \dots & 0 \\
            \vdots & \vdots & \vdots & \ddots & \vdots \\
            0 & \dots & \dots & -1 & 1
        \end{bmatrix}
        \begin{bmatrix}
            u_1 \\
            u_2 - u_1 \\
            u_3 - u_2 \\
            \vdots \\
            u_N - u_{N-1}
        \end{bmatrix}
        =
        \begin{bmatrix}
            u_1 \\
            u_2 - 2u_1 \\
            u_3 - 2u_2 + u_1 \\
            \vdots \\
            u_N - 2u_{N-1} + u_{N-2}
        \end{bmatrix}
    \]
\end{itemize}

\subsection{Scelta del parametro di regolarizzazione $\gamma$ e criterio di discrepanza}
Confronto grafico sul problema simulato di Hunt (\textit{Arbitrary Units}, $t \in [0, 1000]$):
\begin{figure}[h]
    \centering
    \includegraphics[width=1\textwidth]{Immagini/effettisullaregolarizzazione.png}
    \label{fig:7-10}
\end{figure}
\begin{itemize}
    \item Pannelli superiori (\textit{(Under) regularized Deconvolution} vs \textit{Output Noisy Samples \& Reconvolution}): parametro di regolarizzazione troppo piccolo ($\gamma = 5$), la stima presenta ancora oscillazioni spurie.
    \item Pannelli inferiori (\textit{(Over) regularized Deconvolution} vs \textit{Output Noisy Samples \& Reconvolution}): parametro di regolarizzazione troppo grande ($\gamma = 400$), la stima risulta eccessivamente appiattita (\textit{over-smoothing}) e sottostima i picchi.
\end{itemize}

\noindent\textbf{Ruolo del parametro $\gamma$}

Il peso relativo dato ai dati e alla regolarità della soluzione viene gestito dal parametro di regolarizzazione $\gamma$:
\begin{itemize}
    \item Incrementando $\gamma$ si dà maggior peso alla regolarità della soluzione che al fitting dei dati.
    \item Al contrario, riducendo $\gamma$, il costo della regolarizzazione viene ridotto e la fedeltà dei dati diventa importante.
\end{itemize}

La scelta del parametro di regolarizzazione è un problema cruciale:
\begin{itemize}
    \item valori troppo alti portano a stime di $u$ troppo \textit{smooth} che non permettono di spiegare i dati (\textit{over-smoothing});
    \item valori troppo bassi portano a soluzioni mal-condizionate di $u$ che fittano i dati ma con oscillazioni spurie dovute alla loro eccessiva sensibilità al rumore.
\end{itemize}


\noindent Diversi sono gli approcci proposti per la scelta del parametro di regolarizzazione: \textbf{discrepanza}, \textbf{cross-validazione}, \textbf{cross-validazione generalizzata}, \textbf{L-curve}.

\noindent\textbf{Criterio di discrepanza (approccio empirico):}
\begin{itemize}
    \item Confrontiamo il vettore dei residui:
    \[
        r = z - G\hat{u}
    \]
    dove $r$ dipende da $\gamma$ (essendo $\hat{u} = (G^T \Sigma_v^{-1} G + \gamma F^T F)^{-1} G^T \Sigma_v^{-1} z$).
    \item Poiché $v = z - Gu$, allora si può dire che il vettore dei residui è una stima dell'errore di misura:
    \[
        \hat{v} = r \qquad \Longrightarrow \qquad r^T r \cong E[v^T v]
    \]
    \item Poiché $r^T r$ dipende da $\gamma$, bisogna iterativamente aggiustare $\gamma$ fino a che:
    \[
        \sum_{i=1}^{N} r_i^2 \cong \operatorname{trace}(\Sigma_v) = \sum_{i=1}^{N} \sigma_i^2
    \]
\end{itemize}
Vogliamo che il residuo sia simile al rumore di misura, in questo modo il fitting dei dati è soddisfacente e la stima non è eccessivamente \textit{smooth}.
\begin{figure}[h]
    \centering
    \includegraphics[width=0.5\textwidth]{Immagini/esempiocriteriodidiscrepanza.png}
    \label{fig:7-11}
\end{figure}


% ==============================================================================
\section{Campionamento non uniforme e griglia virtuale}
% ==============================================================================

\subsection{Problema del campionamento}
I segnali biologici spesso sono campionati con un \textit{rate} basso e/o non uniforme.
Il metodo di regolarizzazione è basato su un modello discreto derivato assumendo che l'ingresso sconosciuto è costante in ogni intervallo di campionamento, indipendentemente dalla sua lunghezza.
Se il campionamento è lento, si ha una approssimazione povera del segnale ricostruito.
\begin{figure}[h]
    \centering
    \includegraphics[width=0.6\textwidth]{Immagini/problemacampionamento.png}
    \label{fig:7-12}
\end{figure}

\subsection{Formulazione con griglia virtuale}
Le performance deludenti sono dovute al fatto che si assume che il numero di componenti del segnale di input $u$ è uguale al numero $N$ di misure. Per rimuovere questa ipotesi bisogna usare una griglia di campionamento diversa per input e output.

Consideriamo $\Omega_S$ la griglia di campionamento delle misure e $\Omega_v = \{T_1, T_2, \dots, T_k, \dots, T_V\}$ una griglia di campionamento più fine ($V \gg N$, possibilmente uniforme) per l'ingresso. $\Omega_v$ contiene $\Omega_S$ ma, a parte i campioni coincidenti, è arbitraria e non ha una controparte sperimentale (\textbf{griglia virtuale}).
\begin{figure}[h]
    \centering
    \includegraphics[width=0.5\textwidth]{Immagini/grigliadicampionamento.png}
    \label{fig:7-13}
\end{figure}

Assumendo che l'input $u(t)$ è costante in ogni intervallo temporale della griglia virtuale, si ottiene:
\[
    c_V(T_k) = \int_{0}^{T_k} g(T_k - \tau) u(\tau) \, d\tau = \sum_{i=1}^{k} u_i \int_{T_{i-1}}^{T_i} g(T_k - \tau) \, d\tau
\]

In notazione matriciale:
\[
    c_V = G_V u
\]
dove:
\begin{itemize}
    \item $V$ sta per virtuale;
    \item $c_V$ e $u$ sono i vettori $V$-dimensionali ottenuti campionando $c(t)$ e $u(t)$ sulla griglia virtuale;
    \item $G_V$ è una matrice triangolare inferiore $V \times V$.
\end{itemize}

Indicando con $G$ la matrice $N \times V$ ottenuta cancellando da $G_V$ le righe che non corrispondono ai dati campionati in uscita, si ottiene il vettore delle misure:
\[
    z = G u + v
\]
dove:
\begin{itemize}
    \item $v$ è il vettore $N$-dimensionale dell'errore di misura;
    \item $u$ è il vettore $V$-dimensionale dell'input discretizzato sulla griglia virtuale.
\end{itemize}

\noindent\textbf{Griglia virtuale (3/3):} La stima di $\hat{u}$ è data da:
\[
    \hat{u} = (G^T \Sigma_v^{-1} G + \gamma F^T F)^{-1} G^T \Sigma_v^{-1} z
\]
Dove $F$ ha dimensione $V \times V$. La matrice $(G^T \Sigma_v^{-1} G + \gamma F^T F)$ è invertibile perché somma di una matrice semi-definita positiva e una matrice positiva.

\subsection{Esempio numerico di costruzione della matrice $G$}

Consideriamo una griglia di campionamento non uniforme e infrequente ($N = 5$ campioni agli istanti $1, 3, 6, 10, 15$) immersa in una griglia virtuale uniforme ($V = 15$ istanti):
\begin{center}
\begin{tabular}{rccccccccccccccc}
    \textbf{Griglia di campionamento $\Omega_s$:} & $1$ & $?$ & $3$ & $?$ & $?$ & $6$ & $?$ & $?$ & $?$ & $10$ & $?$ & $?$ & $?$ & $?$ & $15$ \\
    \textbf{Griglia virtuale $\Omega_V$:} & $1$ & $2$ & $3$ & $4$ & $5$ & $6$ & $7$ & $8$ & $9$ & $10$ & $11$ & $12$ & $13$ & $14$ & $15$
\end{tabular}
\end{center}

\noindent Nel sistema completo sulla griglia virtuale $c_V = G_V u$ (di dimensione $15 \times 15$), si eliminano le righe corrispondenti agli istanti non misurati ($c(2), c(4), c(5), c(7), c(8), c(9), c(11), c(12), c(13), c(14)$), conservando unicamente le righe $1, 3, 6, 10, 15$:
\[
    \begin{bmatrix}
        c(1) \\
        c(3) \\
        c(6) \\
        c(10) \\
        c(15)
    \end{bmatrix}
    =
    \begin{bmatrix}
        g_0 & 0 & 0 & 0 & 0 & 0 & 0 & 0 & 0 \\
        g_2 & g_1 & g_0 & 0 & 0 & 0 & 0 & 0 & 0 \\
        g_5 & g_4 & g_3 & \dots & g_0 & 0 & 0 & 0 & 0 \\
        g_9 & g_8 & g_7 & \dots & \dots & \dots & g_0 & 0 & 0 \\
        g_{14} & g_{13} & g_{12} & \dots & \dots & \dots & \dots & g_1 & g_0
    \end{bmatrix}
    \begin{bmatrix}
        u_1 \\
        u_2 \\
        u_3 \\
        u_4 \\
        \vdots \\
        \vdots \\
        \vdots \\
        u_{15}
    \end{bmatrix}
\]

\noindent\textbf{Schema dimensionale del modello con griglia virtuale:}
\[
    \underbrace{z}_{5 \times 1} = \underbrace{G}_{5 \times 15} \, \underbrace{u}_{15 \times 1} + \underbrace{v}_{5 \times 1}
\]

\subsection{Casi di studio ed esempi applicativi}

\noindent\textbf{Esempio: Secrezione insulinica durante test IVGTT }
\begin{figure}[h]
    \centering
    \includegraphics[width=1\textwidth]{Immagini/esempiocasiivgtt.png}
    \label{fig:7-14}
\end{figure}
Confronto tra la deconvoluzione regolarizzata su griglia standard (\textit{Regularized Deconvolution (Discrepancy)}, sinistra) e con griglia virtuale (\textit{Regularized Deconvolution (Discrepancy \& Virtual Grid)}, destra), con relativi pannelli inferiori \textit{Output Noisy Samples \& Reconvolution} (ordinate: $\text{pmol/min}$ per la secrezione e $\text{pmol/ml}$ per la concentrazione plasmatica; ascisse: \textit{Minutes} da $0$ a $200$).

\noindent\textbf{Didascalia originale (Figure 9.9):} \textit{IVGTT insulin secretion case study. Left: Regularized deconvolution obtained with a value of the regularization parameter determined by the discrepancy criterion ($\gamma = 0.0012$) from the noisy samples (lower panel, open circles). Also shown is reconvolution (lower panel, continuous line). Right: Same as left by employing the virtual grid ($\gamma = 0.25$) (adapted from Sparacino et al., 2001).}

\noindent\textbf{Esempio 1: Problema simulato di Hunt con griglia virtuale}

Confronto completo tra griglia standard (\textit{Regularized Deconvolution (Discrepancy)}, $\gamma = 105$) e griglia virtuale (\textit{Regularized Deconvolution (Discrepancy \& Virtual Grid)}, $\gamma = 1237455$), che mostra come l'introduzione della griglia virtuale consenta una ricostruzione quasi continua e altamente fedele del profilo vero $u(t)$.

\noindent\textbf{Didascalia originale (Figure 9.6):} \textit{The Hunt simulated problem. Left panels: Regularized deconvolution vs true input (upper, thick vs thin line) obtained with the value of the regularization parameter determined by the discrepancy criterion ($\gamma = 105$) from the noisy samples (lower, open circles). Also shown is reconvolution (lower, continuous line). Right: Same as left by employing the virtual grid ($\gamma = 1237455$).}

\vspace{12pt}
\hrule
\vspace{8pt}

\noindent\textbf{Riferimento finale della lezione }
\begin{itemize}
    \item \textbf{Questionario 1:} \texttt{bit.ly/poll1h1} (accessibile tramite QR code nella slide finale).
\end{itemize}